\chapter{Despliegue del sistema}
\label{Appendix:despliegue}

El sistema está pensado para operar como un conjunto de servicios
independientes: una base de datos PostgreSQL, un backend Java sobre
Apache Tomcat y un frontend estático servido a través de Nginx. Cuando
se quiere compartir una instancia con terceros (por ejemplo, para la
defensa del proyecto o para demostraciones puntuales) se añade un
cuarto componente, un túnel \texttt{cloudflared}, que expone la
aplicación al exterior sin necesidad de abrir puertos en el anfitrión.
Este apéndice resume las decisiones de despliegue y los modos de
arranque previstos. Los pasos exactos, junto con sus
\emph{troubleshooting} habituales, están documentados en el
\texttt{README.md} de la raíz del repositorio\footnote{\url{https://github.com/hibjan/TFG/blob/main/README.md}},
al que se remite para la secuencia operativa concreta.

\section{Modo de desarrollo}

Durante el desarrollo se priorizan dos cosas: el aislamiento entre
componentes y la velocidad de iteración. Por eso solo PostgreSQL se
ejecuta como contenedor, mientras que backend y frontend se levantan
directamente en la máquina del desarrollador. El contenedor de
PostgreSQL se arranca con \texttt{docker-compose.dev.yml}, que monta el
directorio \texttt{database/} sobre el punto de inicialización de la
imagen oficial, lo que provoca que el esquema (descrito más adelante)
se cree automáticamente en el primer arranque sin pasos manuales.

El backend se compila con Maven y se despliega en una instalación
local de Apache Tomcat 10.1 a través de Eclipse, lo que aporta
recompilación incremental y depuración paso a paso. Las credenciales
de la base de datos se inyectan al servidor como variables de entorno,
leyendo los nombres definidos en \texttt{.env.db}
(\texttt{DB\_HOST}, \texttt{DB\_NAME}, \texttt{DB\_USER},
\texttt{DB\_PASS} y \texttt{DB\_PORT}), sin que valor alguno quede
embebido en el código ni en el WAR. El frontend, por su parte, se sirve
con el servidor de desarrollo de Vite, que ofrece recarga en caliente
sobre los cambios y permite apuntar a un backend remoto modificando
únicamente \texttt{VITE\_API\_BASE\_URL} en
\texttt{.env.development.local}.

Esta separación introduce una particularidad: durante el desarrollo,
frontend y backend se sirven en orígenes distintos (puertos diferentes
del mismo \texttt{localhost}), lo que obliga a configurar
explícitamente CORS y a marcar las cookies de sesión como
\texttt{SameSite=None}. La razón está discutida en el cuerpo del
capítulo de desarrollo; el descriptor \texttt{web.xml} y el
\texttt{context.xml} del backend incorporan esta configuración para
que el navegador acepte la cookie de sesión en peticiones
\emph{cross-site}. En el modo compartido, como se verá a continuación,
el problema desaparece porque ambos servicios quedan tras un mismo
origen.

\section{Instancia compartida con túnel Cloudflare}

Para compartir una instancia funcional con un público externo sin
exigirle instalar nada, el stack completo se levanta con un único
comando de \texttt{docker compose} sobre \texttt{docker-compose.prod.yml},
que construye y orquesta los cuatro servicios y los conecta en una red
interna privada. La elección de empaquetar también el backend y el
frontend, y no solo la base de datos, responde a dos motivos. Por un
lado, evita depender del entorno del anfitrión (versiones de Java, de
Node, de Tomcat) y hace el despliegue reproducible. Por otro, permite
unificar el frontend estático y la API tras un mismo origen mediante un
reverse proxy de Nginx, lo que elimina por completo las complicaciones
de CORS y cookies que sí están presentes durante el desarrollo.

El servicio de frontend utiliza una imagen multi-etapa que primero
ejecuta \texttt{npm run build} y a continuación copia los artefactos a
una imagen \texttt{nginx:alpine} con una configuración personalizada.
Esa configuración sirve los archivos estáticos en la raíz y hace de
proxy hacia el contenedor del backend para las rutas bajo
\texttt{/backend/}, de modo que el navegador percibe siempre un único
dominio. El servicio de backend sigue un esquema análogo: una primera
etapa compila el WAR con Maven y una segunda etapa lo despliega sobre
\texttt{tomcat:10.1-jdk17}, aplicando además el ajuste de
\texttt{CookieProcessor} necesario para sesiones \emph{cross-site}.

\begin{figure}[!htb]
  \centering
    \hspace*{2.5cm}%
  \begin{tikzpicture}[
    node distance=1cm and 1cm,
    every node/.style={font=\small},
    box/.style={draw, rectangle, rounded corners, minimum width=5cm,
                minimum height=0.9cm, align=center, fill=orange!10},
    db/.style={draw, cylinder, shape border rotate=90, aspect=0.25,
               minimum width=3cm, minimum height=1.2cm, align=center,
               fill=blue!10},
    arrow/.style={Latex-Latex, thick}
  ]
    \node[box, fill=yellow!40] (user)   {Usuario (navegador)};
    \node[box, fill=red!10,    below=of user]    (cf)    {Cloudflare Edge\\\ttfamily https://*.trycloudflare.com};
    \node[box, fill=red!10,    below=of cf]      (cfd)   {cloudflared (túnel saliente)};
    \node[box,                 below=of cfd]     (nginx) {frontend --- Nginx\\(\texttt{:80} en el contenedor)};
    \node[box,                 below=of nginx]   (tom)   {backend --- Tomcat\\(\texttt{:8080} en el contenedor)};
    \node[db,                  below=of tom]     (pg)    {PostgreSQL\\(\texttt{:5432})};

    \draw[arrow] (user)  -- node[right]{HTTPS} (cf);
    \draw[arrow, pos=0.65] (cf)    -- node[right]{túnel cifrado} (cfd);
    \draw[arrow] (cfd)   -- (nginx);
    \draw[arrow] (nginx) -- node[right]{proxy \texttt{/backend/}} (tom);
    \draw[arrow] (tom)   -- node[right]{JDBC} (pg);

    \node[draw, dashed, fit=(cfd)(nginx)(tom)(pg),
          inner ysep=0.6cm, inner xsep=1.5cm,
          label=right:{Docker Compose}] {};
  \end{tikzpicture}
  \caption{Topología de la instancia compartida con túnel Cloudflare.}
  \label{fig:prodDeployment}
\end{figure}

El componente diferencial es \texttt{cloudflared}, que abre un túnel
saliente hacia la red de Cloudflare y obtiene a cambio una URL pública
en el subdominio \texttt{*.try\allowbreak cloud\allowbreak flare.com}. Esto permite exponer la
aplicación al exterior sin abrir puertos en el cortafuegos del
anfitrión, sin configurar DNS y sin gestionar certificados TLS, ya que
el cifrado y la terminación los proporciona Cloudflare en su extremo
del túnel. La URL concreta se extrae de los logs del contenedor tras
el arranque y puede compartirse directamente con quien deba acceder a
la instancia. La topología resultante, presentada de forma coherente
con la Figura~\ref{fig:arquitectura} del capítulo de arquitectura, se
recoge en la Figura~\ref{fig:prodDeployment}.

\FloatBarrier

\section{Base de datos e ingesta inicial}

\begin{sloppypar}
El esquema relacional, descrito en el cuerpo de la memoria, se
materializa en \texttt{database/\allowbreak 01\_\allowbreak schema.sql}, que crea las
cinco tablas (\texttt{datasets}, \texttt{collections}, \texttt{entities},
\texttt{metadata} y \texttt{reference}) junto con sus claves foráneas e índices.
\end{sloppypar} Este archivo se monta como volumen en el directorio de
inicialización de la imagen oficial de PostgreSQL, lo que provoca que se
ejecute automáticamente la primera vez que el contenedor arranca con un
volumen vacío. La convención de prefijos numéricos en los nombres de
los scripts SQL permite añadir migraciones futuras sin alterar el
mecanismo: PostgreSQL los ejecutará en orden lexicográfico durante el
mismo arranque inicial.

Una vez creado el esquema, la carga de un dataset concreto se realiza
invocando \texttt{populate\_db\_jsonl.py} desde la máquina anfitriona,
según se detalla en el Apéndice~\ref{Appendix:integracion}. El script
publica los datos a través del puerto expuesto por el contenedor de
PostgreSQL, leyendo las credenciales de \texttt{.env.db}. La separación
entre creación del esquema (automática en el primer arranque) y carga
de datasets (manual, posterior) hace el sistema flexible: la base de
datos puede convivir con varios datasets independientes, cargados en
momentos distintos, y un fallo durante la ingesta no compromete la
estructura ya validada.

Si en algún momento se desea regenerar el estado de la base de datos
por completo, basta con detener el contenedor descartando el volumen
asociado, lo que provoca que el siguiente arranque vuelva a ejecutar
\texttt{01\_schema.sql} sobre un volumen vacío. Esta operación elimina
todos los datasets cargados previamente y debe usarse, por tanto, con
cuidado; el \texttt{README.md} la documenta como recurso de
\emph{troubleshooting} para los casos en los que el estado de la base
de datos haya quedado inconsistente.
